\section{Theory}\label{app:theory}
% Propositions 1-4 (PLAN v7 A3). Statements corrected by the numerical checks of the seventh round (unit T); proofs are written by the
% lead and Claude in conversation (PLAN v7 A2/A3) and are pending.
\paragraph{Setting and assumptions.}
An ideal harmonic code places the centred slot encoding of state $x$ in orthogonal planes, $c(x) = \mu + \sum_k a_k
(\cos(2\pi k x/N)\,u_k + \sin(2\pi k x/N)\,v_k)$, with the Nyquist plane $k=N/2$ one-dimensional. Noise is isotropic and
errors are treated to first order. Each statement below says where it matches the data and where it does not.

\paragraph{Proposition 1 (rank threshold).} A residual linear operator that maps every $c(x)$ exactly onto $c(x+1)$ has rank at least
the dimension of the span of the centred code, $N-1$ when every harmonic is present. On ideal codes the minimal exact rank of
closed-form RRR equals this dimension in \CLtheoryPoneHolds{} of \CLtheoryPoneCases{} settings, including codes with missing
harmonics. Exactness is the right notion: under a nearest-state readout a chain can compose at a lower rank when the dropped plane
carries little readout weight. \TODO{proof (lead + Claude).}

\paragraph{Proposition 2 (least squares drops the slowest plane per dimension).} When rank is scarce, reduced-rank regression keeps
the planes with the largest one-step energy per dimension, $4a_k^2\sin^2(\pi k/N)/\dim_k$. Ranking per dimension matches the planes kept
by the fitted operator at \CLtheoryPtwoDim{} of \CLtheoryPtwoRanks{} ranks on ideal codes; ranking per plane, as first stated in the
plan, matches \CLtheoryPtwoPlane{}, all differences coming from the one-dimensional Nyquist plane. If amplitudes do not fall too fast
with $k$, the fundamental is dropped first. On real encodings (exploratory) the most degraded plane of a rank-deficient RRR is the
fundamental in \CLtheoryTtwoAWorst{} of \CLtheoryTtwoARanks{} (backbone, domain, rank) cells. \TODO{proof.}

\paragraph{Proposition 3 (freeze signatures).} Under the weighted nearest-state readout, with the fundamental frozen a chain succeeds
at small displacements and fails near half a turn when the frozen plane's weighted distance dominates; it revives at $t=N$, and
alternating chains keep the outcome of step one. Freezing all other planes fails at step one when their weighted one-step distance
exceeds that of the fundamental. The closed-form penalty equals the vector readout, and revival and alternation hold, in all
\CLtheoryPthreeTrials{} random configurations; with the measured amplitudes and the locked plane weights, the predicted success by
displacement matches the observed frozen-fundamental chains at \CLtheoryTtwoBAgree{} of \CLtheoryTtwoBN{} points (exploratory).
\TODO{proof.}

\paragraph{Proposition 4 (one-cycle cancellation; not confirmed).} If the off-manifold error of a step is linear in the centred code,
the errors of one full one-way cycle sum to zero while alternating paths accumulate them in pairs. Under the linear-error assumption the cancellation is an identity. Of its preregistered tests on held-out templates
(Section~\ref{sec:scope}) only the drift asymmetry was confirmed (alternating chains drift off the manifold far more than one-way chains);
the tests of first-order cancellation and of the inverse operator failed. It remains an unconfirmed
explanation. \TODO{proof.}
